\section{正则化}

\subsection{为什么需要正则化}

损失函数由两部分组成：

$$L(W) = \underbrace{\frac{1}{N}\sum_{i=1}^{N} L_i(f(x_i, W), y_i)}_{\text{Data Loss}} + \underbrace{\lambda R(W)}_{\text{Regularization}}$$

\begin{figure}[H]
\centering
\includegraphics[width=0.95\textwidth]{slides-images/slide-011.jpg}
\caption{总损失 = 数据损失 + 正则化项，$\lambda$ 控制正则化强度\protect\footnotemark}
\end{figure}
\footnotetext{Slides 第15页。}

\textbf{Data Loss} 衡量模型对训练数据的拟合程度——越低越好。\textbf{正则化项} $R(W)$ 的目的则相反：它故意让模型在训练数据上表现稍差，以换取在未见数据上更好的泛化能力。

\begin{figure}[H]
\centering
\includegraphics[width=0.95\textwidth]{slides-images/slide-012.jpg}
\caption{过拟合 vs 泛化：$f_1$ 完美拟合训练数据但过拟合，$f_2$ 拟合不那么紧密但泛化更好\protect\footnotemark}
\end{figure}
\footnotetext{Slides 第16页。}

\begin{knowledgebox}{奥卡姆剃刀与正则化}
正则化的直觉与奥卡姆剃刀（Occam's Razor）一致：当多个假设都能解释数据时，优先选择最简单的那个。复杂模型（如 $f_1$）虽然完美拟合训练数据，但在新数据上表现不佳；而简单模型（如 $f_2$）虽然训练误差略高，但泛化能力更强。
\end{knowledgebox}

\subsection{L2 正则化与 L1 正则化}

\textbf{L2 正则化}（权重衰减）：

$$R(W) = \sum_{k} \sum_{l} W_{k,l}^2$$

对权重矩阵中每个元素取平方后求和。L2 正则化倾向于让权重值分散且均匀较小。

\textbf{L1 正则化}：

$$R(W) = \sum_{k} \sum_{l} |W_{k,l}|$$

对权重矩阵中每个元素取绝对值后求和。L1 正则化倾向于产生稀疏权重（很多值为零）。

\begin{figure}[H]
\centering
\includegraphics[width=0.95\textwidth]{slides-images/slide-013.jpg}
\caption{L2 和 L1 正则化的公式及其对权重分布的不同偏好\protect\footnotemark}
\end{figure}
\footnotetext{Slides 第17页。}

\begin{importantbox}{L2 vs L1 正则化的直觉}
考虑两组权重 $\mathbf{w}_1 = [1, 0, 0, 0]$ 和 $\mathbf{w}_2 = [0.25, 0.25, 0.25, 0.25]$，输入为 $\mathbf{x} = [1, 1, 1, 1]$，两者的点积相同（都等于 1），数据损失一样。但：
\begin{itemize}
    \item L2 偏好 $\mathbf{w}_2$：$0.25^2 \times 4 = 0.25$ vs $1^2 = 1$（$\mathbf{w}_2$ 的 L2 罚项低 4 倍）
    \item L1 对两者无差异：$0.25 \times 4 = 1$ vs $1 = 1$（L1 罚项相同）
\end{itemize}
因此 L2 促进权重分散，L1 促进权重稀疏。
\end{importantbox}

\begin{warningbox}{L1 等价性陷阱}
当被问``L1 正则化偏好 $\mathbf{w}_1$ 还是 $\mathbf{w}_2$？''时，很多人直觉回答 $\mathbf{w}_1$（因为 L1 促进稀疏性）。但实际上在这个例子中，两者的 L1 惩罚值完全相等！L1 在实际优化过程中会产生稀疏解，但这并不意味着任意两组权重的 L1 值不同。
\end{warningbox}

\subsection{正则化的作用总结}

为什么要正则化模型？主要有三个原因：

\begin{enumerate}
    \item \textbf{表达权重偏好}：根据问题特性选择 L1（稀疏）或 L2（分散）
    \item \textbf{简化模型}：防止过拟合，提升测试性能
    \item \textbf{改善优化}：L2 正则化可以使损失函数更凸（$y = x^2$ 是凸函数），有利于优化
\end{enumerate}

\begin{figure}[H]
\centering
\includegraphics[width=0.95\textwidth]{slides-images/slide-015.jpg}
\caption{正则化的三个作用：表达偏好、简化模型、改善优化\protect\footnotemark}
\end{figure}
\footnotetext{Slides 第19页。}

$\lambda$ 是一个超参数，控制正则化的强度：$\lambda = 0$ 表示无正则化，$\lambda$ 越大正则化越强。通常通过验证集来选择最优的 $\lambda$。

\subsection{本章小结}

正则化通过在损失函数中添加权重惩罚项来防止过拟合。L2 正则化偏好分散的小权重，L1 正则化偏好稀疏权重。超参数 $\lambda$ 控制正则化与数据拟合之间的权衡。在后续课程中还会学到更高级的正则化技术（如 Dropout）。

%% ========== 优化基础 ==========

